\scanpage{310}
\mspage{9}
qu'on donne la préférence aux \deleted{\unsure{chaînes}} cochaînes ou aux chaînes.
(Je n'ai pas encore besoin urgent d'une telle convention de notation\ldots)

Dans le cas d'une dualité
\diagram{\begin{tikzcd}[column sep=large]
 \vee:A\arrow[r,shift left=1ex,"\approx"]&B^o\arrow[l,shift left=1ex]
\end{tikzcd}}
on étend les conventions de montée et descente des indices aux objets à
plusieurs indices, ainsi
\[
 ({K_*}^{\cdot})^\vee={K^*}_{\cdot}\deleted{\unsure{\}}}.
\]
Dans la suite des objets mixtes $(\unsure{*|R})$ de la page précédente,
on a écrit l'un au dessus de l'autre des objets qui se correspondent par dualité.

\uline{Remarque} Il est essentiel pour notre système de conventions de définir
(disons) des bicomplexes de cochaînes (objets de $A^{\cdot\cdot}$) avec deux
op. diff. $d_{(1)}^{\deleted{\unsure{*}}\cdot}$ et $d_{(2)}^{\cdot}$ qui
\uline{commutent} --- on n'utilisera donc pas les conventions de
Cartan-Eilenberg. Il faudra donc un \uline{signe} à la diff. totale.

Du point de vue de la prop. 1 et du yoga \unsure{correspondant}
``\unsure{description} d'invariants associés à des $X_*\in\operatorname{Ob}Ss$'',
on peut dire que le \unsure{th.ème|théorème} de Dold-Puppe
\margin{311}
